\chapter*{Introduction}
\addcontentsline{toc}{chapter}{Introduction}

Ultrashort laser processing enables local, controllable modification of metallic surfaces and thin films. Its usefulness in materials engineering depends on predicting when a pulse will heat, melt, or damage a target, and on doing so across materials, film thicknesses, optical conditions, and laser parameters without relying solely on repeated experiments.

In femtosecond laser--metal interactions, absorbed energy is deposited first in the electronic subsystem and is then transferred to the lattice through electron--phonon coupling. This transient nonequilibrium makes the damage threshold a physically informative observable: it compresses a coupled optical and thermal process into a fluence scale that can be compared across simulations, experiments, and material choices.

The parabolic two-temperature model provides the main modelling bridge in this thesis. It combines optical absorption, ballistic and finite-thickness corrections, electron and lattice heat capacities, electron--phonon coupling, and thermal conduction into a tractable description of ultrafast heating. In this role, it is used to compute temperature histories and melting-onset criteria rather than only to state the governing balance laws.

These calculations are numerically demanding because diffusion and electron--phonon exchange introduce stiff time scales. A central part of the thesis is therefore the construction and assessment of stable time-integration schemes for the pTTM, including implicit and implicit--explicit treatments that keep the stiff terms under control while retaining practical runtime for parameter studies.

The same numerical setting supports damage-threshold estimation. Thresholds are obtained by combining pTTM simulations with bracketing and bisection procedures for melting onset, and by comparing the resulting trends across representative metallic systems. These calculations establish the physical and numerical context for later surrogate modelling.

The final layer of the thesis is surrogate-assisted threshold estimation. Rather than learning the threshold directly, the surrogate models learn the maximum lattice temperature \(T_{l,\max}\) over synthetic two-dimensional pTTM datasets; threshold fluences are then obtained by inversion of the predicted temperature--fluence response. This keeps the learned target tied to the thermal model while providing computationally efficient threshold estimates and uncertainty intervals. The main component is the two-dimensional linear pTTM surrogate study, while Chapter~4 also reports a nonlinear extension at the \(\num{20000}\) accepted-sample scale. For Ni and Ti, the nonlinear analysis reduces the discrepancy relative to the linear failure case and reports internally generated nonlinear threshold estimates. Larger-scale nonlinear datasets, expanded nonlinear thermophysical curves for Cr and Al, and multi-target learning remain future work.

Chapter~1 introduces the full physical background for ultrashort laser--metal heating, optical absorption, ballistic and finite-thickness corrections, thermophysical closures, laser heat sources, DLIP laser-interference fields, LIPSS formation, and pTTM model variants. Chapter~2 develops and validates the numerical pTTM schemes. Chapter~3 assembles the optical, thermophysical, DLIP heat-source, laser-interference, LIPSS-relevant source-patterning, and threshold-search components for representative real-material studies. Chapter~4 develops the surrogate-modelling component, reports the linear pTTM surrogate results, and incorporates the nonlinear pTTM extension.
